\section{Background: Circuit Theory}
\subsection{Preliminaries} 
In the context of neural network interpretability, a circuit can be conceptualized as a human-interpretable subgraph that is dedicated to executing specific tasks within a neural network model \cite{olah2020zoom,wang2023interpretability,lv2024interpreting,acdc,bereska2024mechanistic}. 
When we visualize a neural network model as a connected directed acyclic graph (DAG), denoted as $\mathcal{G}$, the individual nodes represent the various components involved in the forward pass, such as neurons, attention heads, and embeddings. 
The edges symbolize the interactions between these components, including residual connections, attention mechanisms, and projections. 
A circuit, represented as $\mathcal{C} \subseteq \mathcal{G}$, emerges as a significant subgraph of $\mathcal{G}$ that is responsible for particular behaviors or functionalities.
In this paper, we focus on the Transformer decoder architecture to conduct our experiments. 
The residual stream of Transformers has been demonstrated to be a valuable tool for mechanistic interpretability in recent works \cite{elhage2021mathematical,yu2024locating}. 
The Transformer architecture typically starts with token embeddings, followed by a sequence of ``residual blocks'' and concludes with a token unembedding. 
Each residual block comprises an attention layer and an MLP layer, both of which ``read'' their input from the residual stream (via a linear projection) and ``write'' their output back to the residual stream through an additive projection.
We can consider an attention head $A_{l,j}$ (the $j$th attention head in layer $l$) as operating on the residual stream from the previous layer, $R_{l-1}$. 
Given that $R_0=I$ (where $I$ represents the input embeddings), we can reinterpret attention head $A_{l,j}$ as processing the cumulative output of all previous attention heads and MLPs and input embedding, treating each node in the previous layers as separate input arguments. 
Similarly, an MLP node $M_l$ can be seen as operating on the cumulative output of all previous attention heads and MLPs and input embedding, and the output node $O$ operates on the sum of the input embeddings and the outputs of all attention heads and MLPs.
The following equations represent the residual connections in the Transformer model, where $R_l$ is the residual stream at layer $l$, and $\text{Input}^A_l$ and $\text{Input}^M_l$ are the inputs to the attention and MLP layers, respectively:
\begin{align*}
R_l = R_{l-1}+\sum_j A_{l,j}+M_l, R_{0} = I \\
\text{Input}^A_l = I+\sum_{l'<l} \left(M_{l'}+\sum_{j'} A_{l',j'}\right) \\
\text{Input}^M_l=I+\sum_{l'<l}M_{i'} + \sum_{l'\leq i}\sum_{j'}A_{l',j'}
\end{align*}
The computational graph $\mathcal{G}$ of the Transformer represents the interactions between attention heads and MLPs. 
The nodes in $\mathcal{G}$ encompass the input embedding $I$, attention heads $A_{l,j}$, MLPs $M_l$, and the output node $O$, denoted as $N = \{I, A_{l,j}, M_{l}, O\}$.
The edges in the model represent the connections between these nodes, $E = \{(n_{x}, n_{y}), n_{x}, n_{y} \in N\}$.
A circuit $\mathcal{C}$ is meticulously constructed to govern specific behaviors within the model, comprising a selection of nodes $N_{\mathcal{C}}$ and edges $E_{\mathcal{C}}$ that are critical to the successful execution of the tasks at hand, expressed as $\mathcal{C} = <N_{\mathcal{C}}, E_{\mathcal{C}}>$.
% Let's denote the OV matrix as $\mathbf{W}_{O V}=\mathbf{W}_O \mathbf{W}_V$ and the $\mathrm{QK}$ matrix as $\mathbf{W}_{Q K}=\mathbf{W}_Q^{\top} \mathbf{W}_K$, for clarity. The parameters represented by $\mathbf{W}_{O V}$ generate outputs corresponding to inputs; The parameters represented by $\mathbf{W}_{Q K}$ generate attention patterns, determining which token's information is moved from and to.
\subsection{Circuit Discovery}
To identify circuits within a language model, a key approach is to examine the model's casual mediation by systematically altering the model's edges and nodes to observe the effects on performance \cite{acdc,pearl2022direct,vig2020investigating}. 
The underlying principle is that critical edges or nodes are those whose removal results in a notable decline in the model's predictive capabilities.
Since the edges in the model's computational graph represent the dependencies between nodes, we can simulate the absence of a particular node-to-node dependency by ablating an edge in the graph.
For example, ablating an edge from $A_{i',j'}$ to $A_{i,j}$ involves replacing the contribution of $A_{i',j'}$ in the input to attention head $A_{i,j}$ with zero (in the case of zero ablation) or with the mean value of head $A_{i',j'}$ (in the case of mean ablation).
The process of identifying critical edges or nodes through ablation can be broken down into the following steps:
i) Overwrite the value of the edge $(n_{x},n_{y})$ with a corrupted value (either zero or mean ablation),
ii) Perform a forward pass through the model with the altered graph,
iii) Compare the output values of the modified model with those of the original model using a chosen metric $S$ (Details in Eq.\ref{metric_S} ).
% If the performance change lower than the threshold $\tau$, we can remove the edge and obtain a new subgraph $\mathcal{G}/{(n_{x},n_{y})}$.
If the performance change is below a predefined threshold $\tau$, we can consider the edge non-critical and remove it to obtain a new subgraph $\mathcal{G}/{(n_{x},n_{y})}$.
In addition to ablation-based methods, recent works have also explored the use of sparse auto-encoders \cite{he2024dictionary,cunningham2023sparse} to identify circuits within language models. 
This approach involves training an auto-encoder to learn a sparse representation of the model's internal structure, which can help reveal the underlying circuitry responsible for specific behaviors or functionalities.


